% GENERATED by build_do_now_tex.py on 2026-04-24 --- do not hand-edit; regenerate instead.
\documentclass[11pt]{article}
\usepackage{preamble}

\begin{document}

\packetheading{Lesson 4-3 \textperiodcentered\ Period 3 \textperiodcentered\ Do Now}{L43 P3}

\vspace{0.25em}
\noindent\textbf{Name:}\ \rule{2in}{0.4pt}\quad
\textbf{Period:}\ \rule{0.5in}{0.4pt}\quad
\textbf{Date:}\ \rule{1in}{0.4pt}
\vspace{0.5em}

\bankitem{Do Now}{%
Reason Explain why (4x^2-7)/(4x^2-7)=1 is a valid identity under the domain of all real numbers except ±\frac{√(7)}{2}.
\answer{Using the difference of two squares formula, (4x^2-7)/(4x^2-7)=\frac{(2x+√(7))(2x-√(7))}{(2x+√(7))(2x-√(7))}. Both factors on the numerator and denominator divide to 1, so (4x^2-7)/(4x^2-7)=1 over all elements of the domain, which is all real numbers except ±\frac{√(7)}{2}.}
}{}
\vspace{2.5in}

\end{document}
